\chapter{难度感知的自适应采样方法}\label{chap:method}

\section{问题定义}

设 prompt 分布为 $x\sim d_0$，策略模型为 $\pi_{\theta}(y\mid x)$，回答 $y$ 经验证器判定后得到奖励 $r(x,y)\in\{0,1\}$。强化学习后训练的目标可以写为
\begin{equation}
    J(\theta)=\mathbb{E}_{x\sim d_0,\,y\sim\pi_{\theta}(\cdot\mid x)}[r(x,y)].
\end{equation}

对给定 prompt $x$，记单次采样答对的真实概率为 $p_x$。若采用固定组大小 $n$ 的组相对更新，组内全对或全错都会导致奖励方差趋于零。相应地，prompt 级“信号坍塌”概率可以表示为
\begin{equation}
    P_{\mathrm{collapse}}(x;n)=p_x^n+(1-p_x)^n.
\end{equation}

对全体训练分布求期望后，可得固定组大小下的总体坍塌率
\begin{equation}
    P_{\mathrm{collapse}}(n)=\mathbb{E}_{x\sim d_0}[p_x^n+(1-p_x)^n].
\end{equation}

该式揭示了一个关键事实：即使所有 prompt 都满足 $0<p_x<1$，只要 $n$ 足够小，整体训练中仍会频繁出现统计意义上的“零信号”事件。因此，本文关心的问题不是简单地判断某题是否可学，而是讨论在预算受限条件下，如何为不同难度 prompt 分配更合适的采样次数。

\section{难度异质性与预算效率分析}

若目标是尽快观察到正负样本差异，则 prompt 的学习价值不仅取决于正确率本身，也取决于少数类样本的获取成本。对二值奖励而言，单样本方差信息可由 $p_x(1-p_x)$ 表征。若进一步考虑在训练后期仍需保留错误样本，则其单位成本信息效率可近似写为
\begin{equation}
    \eta(p_x)\propto p_x(1-p_x)^2.
\end{equation}

当 $p_x=0.9$ 时，$\eta(p_x)=0.009$；当 $p_x=0.3$ 时，$\eta(p_x)=0.147$。这意味着在相同预算下，中等偏难样本能够产生远高于简单样本的有效对比信息。若训练过程继续对高通过率 prompt 强行追求“平衡”的正负样本数量，就会在后期为少量 hard negative 付出不成比例的重采样开销。

基于这一分析，本文主张将退出规则与批次预算设计成难度敏感的形式：对高成功率 prompt，允许更早停止并将预算腾挪给中等难度样本；对低成功率或中等成功率 prompt，则保留额外采样空间，以便更稳定地暴露正负差异。

为了进一步把“预算效率”写成可比较的量，下面考虑固定组大小下首次观察到有效差异的等待时间。若采用平衡型更新，则一个 prompt 在单轮中形成有效信号，等价于组内至少同时出现一个正确样本和一个错误样本。对独立采样近似而言，其单轮成功概率为
\begin{equation}
    q_{\mathrm{bal}}(p_x;n)=1-p_x^n-(1-p_x)^n.
\end{equation}
于是，首次观察到有效信号所需的期望轮数可写为
\begin{equation}
    \mathbb{E}[T_{\mathrm{bal}}]=\frac{1}{q_{\mathrm{bal}}(p_x;n)}.
\end{equation}
该式说明，固定预算的真实低效并不只表现为“某一轮恰好坍塌”，更表现为大量 prompt 需要跨越多个训练轮次，才能偶然获得一次非零梯度机会。当 $p_x$ 很小或很大时，$q_{\mathrm{bal}}$ 都会迅速下降，从而导致等待时间急剧增长。

进一步地，当关注的是训练后期的简单题，即 $p_x\rightarrow 1$ 时，可令 $\delta_x=1-p_x$，则
\begin{equation}
    q_{\mathrm{bal}}(p_x;n)=1-(1-\delta_x)^n-\delta_x^n\approx n\delta_x,
\end{equation}
其中近似利用了 $\delta_x\ll 1$ 时的一阶展开。这意味着，对高通过率 prompt 而言，平衡规则的额外成本本质上由“找到一个错误答案”的等待时间主导，其期望追加开销近似与 $1/(n\delta_x)$ 同阶增长。换言之，当错误样本已经成为少数类时，继续要求平衡退出，会把越来越多预算消耗在寻找稀有 hard negative 上。
\section{几何块自适应采样框架}

为兼顾统计效率与系统效率，本文采用几何块自适应采样框架。对每个 prompt，不再一次性采样固定 $n$ 条回答，而是按块逐轮追加采样。设块序列为
\begin{equation}
    \mathcal{B}=\{2,2,4,8,16\},
\end{equation}
其总预算上限为 $32$。第 $t$ 轮结束后，根据当前累计样本集合 $S_x^{(t)}$ 判断是否继续采样。

该设计具有三点动机。第一，小块起步可以在训练早期快速确认极易题和极难题是否值得继续采样。第二，几何扩张使困难 prompt 获得越来越大的后续预算，而简单 prompt 往往在前两轮即可退出。第三，更规整的块结构更适合与前缀复用和 KV Cache 管理结合，从而减少重复编码开销。

图~\ref{fig:adaptive_pipeline} 给出了本文方法的概念流程图。与固定组采样“一次性生成、一次性更新”的流程不同，本文方法在 prompt 级维护一个动态活跃集合：每一轮只对仍未满足退出条件的样本继续追加预算，其余样本则提前退出。最终，活跃集合与样本池状态共同决定哪些 prompt 会进入后续更新构造阶段。

若与引言中的总览图对照，图~\ref{fig:adaptive_pipeline} 对应的正是其中“方法框架”那一层：总览图说明本文要回答哪些研究问题，而这里的流程图则把这些问题具体落实为活跃集合、几何块预算、退出规则和更新构造四个可操作部件。也正因为如此，相关工作矩阵、总览图和方法流程图在本文中并不是重复表达，而是分别承担“研究定位—全局框架—局部实现”三种不同粒度的叙事功能。

\begin{figure}[!htbp]
    \centering
    \begin{tikzpicture}[>=latex, every node/.style={font=\small}]
        \tikzstyle{box}=[draw, rounded corners, align=center, minimum height=0.95cm, text width=3.0cm]
        \tikzstyle{widebox}=[draw, rounded corners, align=center, minimum height=0.95cm, text width=3.8cm]
        \tikzstyle{line}=[draw, ->, thick]

        \node[box] (prompt) at (0,0) {Prompt 批次\\$\mathcal{X}$};
        \node[widebox] (active) at (4.6,0) {活跃集合初始化\\$\mathcal{A}\leftarrow\mathcal{X}$};
        \node[widebox] (block) at (9.8,0) {按几何块采样\\$2\rightarrow 2\rightarrow 4\rightarrow 8\rightarrow 16$};
        \node[widebox] (pool) at (9.8,-2.2) {累计样本池\\$S_x\leftarrow S_x\cup\{y_i\}$};
        \node[widebox] (rule) at (4.6,-2.2) {退出规则判断\\平衡 / 正样本优先};
        \node[box] (next) at (0,-2.2) {更新活跃集合\\$\mathcal{A}_{next}$};
        \node[widebox] (update) at (4.6,-4.4) {形成有效信号后\\进入下采样与更新构造};

        \path[line] (prompt) -- (active);
        \path[line] (active) -- (block);
        \path[line] (block) -- (pool);
        \path[line] (pool) -- (rule);
        \path[line] (rule) -- (next);
        \draw[->, thick] (next.north) |- (block.west);
        \path[line] (rule) -- (update);
    \end{tikzpicture}
    \bicaption{难度感知自适应采样的概念流程图}{Conceptual pipeline of difficulty-aware adaptive sampling.}
    \label{fig:adaptive_pipeline}
\end{figure}

从实现角度看，这一框架还具有一个重要优点：它不要求训练系统在一开始为每个 prompt 分配固定长度的回答队列，而是允许在轮次边界上以 prompt 为粒度动态调整活跃集合。于是，真正消耗大块预算的只会是那些在前几轮尚未形成可分辨信号的样本，而非全体样本。这种“逐轮淘汰活跃 prompt”的机制，使得自适应采样天然适合与在线训练流水线结合。

\section{两类退出规则}

结合现有实验数据，本文重点分析两类退出规则。

\subsection{平衡退出规则}

平衡退出规则要求当前样本集合中同时出现正确与错误回答，才能认定该 prompt 已经形成有效训练信号。其优势在于显式保留正负样本对照，缺点是在训练后期可能为简单题反复追加采样，只为“寻找一个错误答案”。

若总预算上限为 $m$，则平衡退出在预算耗尽前成功停止的概率可写为
\begin{equation}
    P_{\mathrm{stop}}^{\mathrm{bal}}(p_x;m)=1-p_x^m-(1-p_x)^m.
\end{equation}
相应地，失败事件恰好就是 signal collapse。该表达式说明，平衡规则虽然最符合组相对更新的原始直觉，但它对简单题和极难题都不友好：前者缺少负样本，后者缺少正样本，二者都会把 prompt 长时间保留在活跃集合中。
\subsection{正样本优先退出规则}

正样本优先退出规则在观测到至少一个正确回答后即可停止采样，仅将“预算耗尽仍全错”的 prompt 视作缺失样本。该规则更加节省预算，尤其适用于模型已经具备一定基础能力时的中后期训练阶段。它的核心思想不是否认负样本的重要性，而是承认在有限预算下，简单题上的额外找错成本未必比把预算转移给难题更划算。

在同样的总预算上限 $m$ 下，正样本优先规则的成功停止概率为
\begin{equation}
    P_{\mathrm{stop}}^{\mathrm{pos}}(p_x;m)=1-(1-p_x)^m.
\end{equation}
当 $p_x>0$ 时，总有 $P_{\mathrm{stop}}^{\mathrm{pos}}(p_x;m)\geq P_{\mathrm{stop}}^{\mathrm{bal}}(p_x;m)$，且不等式在 $p_x\in(0,1)$ 上通常严格成立。原因在于，正样本优先只要求观察到“至少一个正样本”，而不再额外要求同时命中稀有负样本。因此，对训练中后期的高通过率 prompt 而言，该规则能以更高概率在较小预算内提前退出，并把节省下来的预算转移给那些仍处于“正样本稀缺”区间的题目。

若进一步用几何分布近似正样本首次出现的等待时间，则其期望采样步数满足
\begin{equation}
    \mathbb{E}[\tau_{\mathrm{pos}}]\approx \frac{1}{p_x}.
\end{equation}
这一定量关系揭示了正样本优先规则的另一个优势：它天然更关注“正样本何时首次暴露”，因此更适合服务于能力显化型训练阶段；而平衡规则则额外叠加了对稀有反例的搜索开销。本文实验中观察到的成本差异，本质上正对应了这两种等待时间结构的不同。

\section{退出规则的停止时间分析}

为了更精确地解释两类退出规则为何会产生不同的预算效率，本文进一步将 prompt 级采样过程写成一个离散停止时间问题。设块序列为 $\mathcal{B}=\{B_1,B_2,\ldots,B_T\}$，累计预算记为
\begin{equation}
    M_t=\sum_{i=1}^{t}B_i,\qquad M_0=0.
\end{equation}

对任意 prompt $x$，若第 $t$ 轮结束后仍未满足退出条件，则称其在第 $t$ 轮后仍处于“存活”状态。记该存活概率为 $q_t(x)$，则期望采样成本可以统一写为
\begin{equation}
    \mathbb{E}[C\mid x]=\sum_{t=1}^{T} B_t\, q_{t-1}(x),\qquad q_0(x)=1.
\end{equation}

这一定义的含义非常直接：第 $t$ 块预算只有在 prompt 经历前 $t-1$ 轮后仍未退出时才会被实际消耗，因此各块大小应与“仍然值得继续投入的概率”共同决定总成本。

\subsection{平衡退出规则的停止时间}

在平衡退出规则下，prompt 之所以会继续留在活跃集合中，仅仅因为当前已观察到的回答仍然全对或全错。于是，在累计预算为 $M_t$ 时，其存活概率满足
\begin{equation}
    q_t^{\mathrm{bal}}(x)=p_x^{M_t}+(1-p_x)^{M_t}.
\end{equation}

因此，平衡规则下的期望成本可写为
\begin{equation}
    \mathbb{E}[C_{\mathrm{bal}}\mid x]=\sum_{t=1}^{T}B_t\Bigl(p_x^{M_{t-1}}+(1-p_x)^{M_{t-1}}\Bigr).
\end{equation}

该式揭示出一个关键性质：当 $p_x$ 接近 $0$ 或 $1$ 时，$q_t^{\mathrm{bal}}(x)$ 会衰减得非常慢，因为 prompt 很容易持续停留在“全错”或“全对”的单侧状态。于是，平衡规则会把大量后续预算投入到那些最不容易形成双侧对照的样本上，这正是其在训练中后期预算效率下降的根源。

\subsection{正样本优先规则的停止时间}

在正样本优先规则下，prompt 继续存活只意味着截至当前仍未观察到任何一个正确回答。于是有
\begin{equation}
    q_t^{\mathrm{pos}}(x)=(1-p_x)^{M_t},
\end{equation}
从而
\begin{equation}
    \mathbb{E}[C_{\mathrm{pos}}\mid x]=\sum_{t=1}^{T}B_t(1-p_x)^{M_{t-1}}.
\end{equation}

与平衡规则相比，正样本优先规则完全移除了 $p_x\rightarrow 1$ 这一端的“找错成本”。当 prompt 已经较容易答对时，它在前几轮就会快速退出；真正继续消耗预算的，主要是那些 $p_x$ 较低、仍有必要继续探索的样本。因此，这一规则天然对应一种更强的预算右移机制：越难获得正样本的 prompt，越能保留后续预算；越容易获得正样本的 prompt，越早退出。

\subsection{单位预算的边际信号收益}

进一步地，可将第 $t$ 轮预算带来的“新增可辨别概率”写为
\begin{equation}
    \Delta_t(x)=q_{t-1}(x)-q_t(x),
\end{equation}
并定义第 $t$ 块的单位预算边际收益为
\begin{equation}
    g_t(x)=\frac{\Delta_t(x)}{B_t}.
\end{equation}

如果 $g_t(x)$ 很低，说明继续为该 prompt 投入这一块预算，能够显著改变训练信号状态的概率并不高。对于平衡规则，这种情况会同时出现在极难样本和极易样本两端；而对正样本优先规则，它主要集中在极难样本一侧。也就是说，后者把“无效追加预算”的区域从两端压缩为一端，因此在平均成本相近时往往具有更好的总体效率。

这一分析也解释了为什么本文采用几何块而不是等长小块：当后续块大小逐轮增大时，系统实际上是在把更大的预算块留给那些在前几轮仍然保持较高 $q_t(x)$ 的样本。换言之，几何块策略并不是随意的工程折衷，而是把“先用廉价小块快速排除明显样本，再对难分样本集中加码”的预算哲学直接编码进采样流程。

\section{训练时采样流程的算法化描述}

为了将上述框架落到可执行流程，算法~\ref{alg:adaptive_sampling} 给出了训练阶段的 prompt 级自适应采样过程。它强调两个关键状态：其一是活跃 prompt 集合，即当前仍需继续采样的题目；其二是每个 prompt 已累计的样本池，用于判断是否满足退出条件。相比“一次性采完再过滤”的被动流程，该算法在每一轮之后都会收缩活跃集合，因此能够将后续预算逐渐集中到最难分辨的样本上。

\begin{algorithm}[!htbp]
    \caption{几何块自适应采样流程}
    \label{alg:adaptive_sampling}
    \small
    \begin{algorithmic}[1]
        \Require Prompt 批次 $\mathcal{X}$，块序列 $\mathcal{B}=\{2,2,4,8,16\}$，退出规则 $\mathcal{R}$
        \Ensure 每个 prompt 的样本池 $S_x$ 及最终是否缺失有效信号的标记
        \State 初始化活跃集合 $\mathcal{A}\leftarrow \mathcal{X}$
        \ForAll{$x\in\mathcal{X}$}
            \State $S_x\leftarrow \varnothing$
        \EndFor
        \For{$b\in\mathcal{B}$}
            \ForAll{$x\in\mathcal{A}$}
                \State 从 $\pi_{\theta}(\cdot\mid x)$ 采样 $b$ 个回答并加入 $S_x$
            \EndFor
            \State $\mathcal{A}_{\mathrm{next}}\leftarrow \varnothing$
            \ForAll{$x\in\mathcal{A}$}
                \If{$\mathcal{R}(S_x)=\mathrm{continue}$}
                    \State 将 $x$ 加入 $\mathcal{A}_{\mathrm{next}}$
                \EndIf
            \EndFor
            \State $\mathcal{A}\leftarrow \mathcal{A}_{\mathrm{next}}$
            \If{$\mathcal{A}=\varnothing$}
                \State \textbf{break}
            \EndIf
        \EndFor
        \ForAll{$x\in\mathcal{X}$}
            \State 根据 $S_x$ 判断是否形成有效信号，并执行后续下采样与更新构造
        \EndFor
    \end{algorithmic}
\end{algorithm}

\section{复杂度与系统实现讨论}

若设初始 prompt 数量为 $N$，第 $t$ 轮仍然活跃的 prompt 数量为 $N_t$，对应块大小为 $B_t$，则总采样成本可以写为
\begin{equation}
    C_{\mathrm{total}}=\sum_t N_t B_t.
\end{equation}

固定组大小策略的成本则为 $C_{\mathrm{fixed}}=N\cdot n$。两者的差异不在于是否都需要为每个 prompt 生成回答，而在于自适应采样允许 $N_t$ 随轮次快速下降。只要简单题与极难题能在前几轮被识别出来，就有机会在后续大块预算到来前从活跃集合中退出，从而压低总体成本。

从显存与缓存角度看，几何块策略还带来另一层收益。对于保留在活跃集合中的 prompt，其前缀往往已经被多轮复用，因此后续继续生成更长回答或更多回答时，能够更自然地复用已有的 KV Cache 状态。换言之，块结构越规整，越容易在系统层面实现“先筛查、后集中扩展”的执行模式；这正是本文强调几何块策略不仅有统计意义，也有工程意义的原因。

\section{硬件感知的指数增长批次采样策略}

前文给出的几何块采样动机主要来自统计效率；但在真实训练系统中，批次结构还会直接影响前缀复用与调度开销。对同一 prompt 的多回答生成而言，推理过程通常由 Prefill 与 Decode 两阶段组成：前者负责对输入前缀编码并构建 KV Cache，后者逐 token 扩展回答。若额外采样能够尽量复用已有前缀状态，则同样的统计预算可以转化为更低的 wall-clock 开销。

为量化这一点，记某个 prompt 的总采样数为 $N$，实际触发 Prefill 的轮次数为 $R$。则可用前缀缓存命中率
\begin{equation}
    \rho = 1 - \frac{R}{N}
\end{equation}
刻画多轮采样带来的额外前缀重算比例。对固定预算一次性生成而言，$R=1$，命中率最高；但这种方式会在大量简单题上执行明显过量的 Decode。相对地，自适应采样的关键不是单独追求最高 $\rho$，而是在保持较高前缀复用的同时，让简单题尽早退出、把额外样本深度留给真正困难的 prompt。

基于这一考虑，本文将几何块结构进一步写成一个硬件感知的指数增长批次策略。设第 $t$ 轮结束时的累计样本数为
\begin{equation}
    N_t = M_0\cdot 2^{t-1},
\end{equation}
其中 $M_0$ 为初始批次大小。于是第 $t$ 轮新增批次满足
\begin{equation}
    B_1=M_0,\qquad B_t=M_0\cdot 2^{t-2}\quad (t\ge 2).
\end{equation}
当 $M_0=2$ 时，批次序列即为 $2,2,4,8,16$。这一结构的优点在于：前两轮仍保持很小的筛查成本，而一旦 prompt 进入后期轮次，后续更大的 batch 可以更充分地分摊单次 Prefill 的固定开销，从而把统计上的“继续探索难题”更自然地映射为系统上的“集中扩展活跃样本”。

若某个 prompt 在第 $k$ 轮停止，则对应的缓存命中率可写为
\begin{equation}
    \rho_{\mathrm{geo}}(k)=1-\frac{k}{N_k}.
\end{equation}
这一定义揭示了几何块策略的两个阶段性特征：简单题通常在很小的 $N_k$ 下快速退出，因此节省了大量 Decode；困难题虽然经历更多轮次，但其后期批次更大，能够更有效地摊销 Prefill 成本。换言之，几何块并不是把固定预算机械拆分成若干轮，而是把“先廉价筛查、后集中扩展”的系统执行逻辑直接编码进采样器中。

\section{动态温度缩放与重要性采样校正}

当某个 prompt 在标准温度下经过多轮采样仍未出现正样本时，一个自然推断是：正确路径可能已经存在于策略分布中，但其概率质量过低，因而难以被固定温度的浅层探索触达。为此，本文进一步将采样温度设计为随累计采样深度逐步增长：
\begin{equation}
    T_t = T_0\cdot\bigl(1+\alpha\log_2(N_t/M_0)\bigr),
\end{equation}
其中 $T_0$ 为基准温度，$\alpha$ 为缩放系数。该设计使浅层采样阶段保持接近标准分布，而深层采样阶段逐步提高探索强度，更适合暴露低概率但潜在正确的推理路径。

但一旦跨轮次温度不同，不同样本实际上来自不同的采样分布。若仍直接把这些样本混合用于同一目标策略更新，就会引入分布偏差。为控制这一问题，本文采用重要性权重
\begin{equation}
    \rho_i=\frac{\pi_{\theta}(y_i\mid x;T_{\mathrm{target}})}{\pi_{\theta}(y_i\mid x;T_{t_i})}
\end{equation}
对多温度样本进行校正，并使用截断版本 $\rho_i^{\mathrm{clip}}=\min(\rho_i,\rho_{\max})$ 抑制极端权重带来的方差膨胀。这里 $T_{t_i}$ 表示样本 $y_i$ 所在轮次的实际采样温度，$T_{\mathrm{target}}$ 则表示训练更新所对应的目标策略温度。

在此基础上，可进一步用自归一化重要性采样估计 prompt 级基准线：
\begin{equation}
    b_{\mathrm{IS}}(x)=\frac{\sum_i \rho_i\, r_i}{\sum_i \rho_i},
\end{equation}
并据此构造经校正的优势函数
\begin{equation}
    A_{\mathrm{IS}}(x,y_i)=r_i-b_{\mathrm{IS}}(x).
\end{equation}
于是，一个统一的加权目标可以写为
\begin{equation}
    L(\theta)=-\sum_i \rho_i^{\mathrm{clip}}\,A_{\mathrm{IS}}(x,y_i)\,\log\pi_{\theta}(y_i\mid x;T_{\mathrm{target}}).
\end{equation}
它的含义是：深层采样阶段允许策略以更高温度进行探索，但更新阶段仍通过重要性权重把不同温度下观察到的信号拉回同一目标分布，避免训练过程过度偏向偶然采中的高温样本。

因此，本文的方法框架可以被理解为三个部件的协同：几何块负责把预算按难度层次展开，正样本优先退出负责把简单题尽早移出活跃集合，而动态温度与重要性采样校正则负责让后期深度采样既更有探索性，又尽量保持统计一致性。它们共同构成了本文所强调的“统计效率—系统效率”联合设计视角。

\section{本章小结}

本章从概率建模、预算效率和系统组织三个角度给出了难度感知自适应采样的形式化描述。核心思想并不复杂：在训练过程中，不同 prompt 对额外采样预算的边际收益并不相同，因此采样次数不应是固定常数，而应由样本难度、退出规则以及批次组织方式共同决定。围绕这一思想，本文进一步提出了几何块调度、活跃集合迭代、硬件感知的指数增长批次策略，以及与之配套的动态温度缩放和重要性采样校正框架，为下一章的实验比较提供了更完整的实现对象与分析基线。
